% Direct-authority diplomatic French for printed p.205.
\clearpage
\SourcePage{205}

\LayerTarget{tohoku:proof:corollary-5.3.1:2:p1}
\noindent On a $\mathrm I_2^{p,q}=0\ \mathrm{mod}\ C(P)$ pour $q>0$, car en
vertu du théorème 5.3.1 $\mathbf H^q(G,A)$ est un faisceau annulé par $n$,
donc $H^p(X,\mathbf H^q(G,A))$ est annulé par $n$ pour tout $p$, donc est nul
mod. $C(P)$. Le corollaire 2 en résulte aussitôt.

\LayerTarget{tohoku:application:fixed-point-free-case-5.3}
Si les $n_x$ sont tous réduits à 1, c'est à dire si tout $g\in G$, $g\ne e$,
n'a pas de point fixe (on dit alors simplement que $G$ \textit{opère sans
points fixes}) on trouve $\mathbf H^n(G,A)=0$ pour $n>0$, donc
$H^*(X,G,A)=H^*(Y,A^G)$, ce qui est aussi évident a priori, à cause de
l'isomorphisme signalé à la fin de 5.1 entre la catégorie des faisceaux sur
$Y$ et la catégorie des $G$-faisceaux sur $X$, quand $G$ est un groupe discret
d'homéomorphismes sans points fixes. On trouve donc dans ce cas l'énoncé
classique :

\LayerTarget{tohoku:corollary:5.3.1:3}
\noindent\textsc{Corollaire 3.} \textit{Si $G$ opère ``sans points fixes'',
alors $H^*(Y,A^G)$ est l'aboutissement d'une suite spectrale cohomologique dont
le terme initial est $H^p(G,H^q(X,A))$.}

\LayerTarget{tohoku:remark:spectral-sequence-construction-5.3}
\noindent\textsc{Remarque.} Par définition, la deuxième suite spectrale du
théorème 5.2.1. s'obtient en prenant une résolution $C(A)$ de $A$ par des
faisceaux $\Gamma_X$-acycliques, et en prenant la deuxième suite spectrale du
foncteur $\Gamma^G$ par rapport au complexe $\Gamma_X(C(A))$. En général, la
première suite spectrale de $\Gamma^G$ par rapport à ce complexe ne s'identifie
évidemment pas à la première suite spectrale du théorème 5.2.1 (prendre une
résolution injective de $A$ !). Il en est cependant bien ainsi si on prend pour
$C(A)$ la \textit{résolution canonique de $A$} (cf. 3.3), comme on voit par un
raisonnement analogue à celui de la proposition 4.3.2. Il en est encore ainsi
si $X$ et $Y$ sont paracompacts et si on prend pour $C(A)$ une résolution de
Cartan de $A$, i.e. $C(A)=A\otimes_{\mathbf Z}C$, où $C$ est un ``faisceau
fondamental'' au sens de Cartan-Leray, comme on peut montrer par des
raisonnements assez différents (nous n'aurons pas à nous servir de ce fait).
Par suite, la suite spectrale I est bien la suite spectrale de [10] dont le
terme initial avait été laissé non explicité.

\LayerTarget{tohoku:section:5.4}
\bigskip
\noindent\textbf{5.4. Transformation de la première suite spectrale.} Nous
supposons toujours la condition (D) du numéro précédent satisfaite (bien que
ce ne soit pas absolument essentiel).

\LayerTarget{tohoku:setup:transformation-first-spectral-sequence}
Soit $Y_0$ une partie fermée de $Y$ contenant les supports des
$\mathbf H^n(G,A)$ pour $n>0$ (il suffit par exemple, en vertu du théorème
5.3.1, que $Y_0$ contienne l'ensemble des $y\in Y$ tels que $G_y\ne e$). Soit
$X_0=f^{-1}(Y_0)$, soit $V$ le complémentaire de $Y_0$ et $U$ le
complémentaire de $X_0$. Considérons les suites exactes :

\LayerTarget{tohoku:sequence:5.4-support-exact-sequences}
\[
\begin{gathered}
0\longrightarrow A_U\longrightarrow A\longrightarrow A_{X_0}\longrightarrow0,\\
0\longrightarrow(A^G)_V\longrightarrow A^G\longrightarrow
  (A^G)_{Y_0}\longrightarrow0.
\end{gathered}
\]
Comme on a manifestement $(A_U)^G=(A^G)_V$, on obtient un diagramme
commutatif de suites de cohomologie

\LayerTarget{tohoku:diagram:5.4.1}
\[
\tag{5.4.1}
\begin{array}{c@{\;}c@{\;}c@{\;}c@{\;}c@{\;}c@{\;}c}
\mathinner{\cdotp\cdotp\cdotp\cdotp}H^n(Y,(A^G)_V)&\to&H^n(Y,A^G)&\to&H^n(Y,(A^G)_{Y_0})&\to&
H^{n+1}(Y,(A^G)_V)\!\to\!\mathinner{\cdotp\cdotp\cdotp\cdotp}\\[-0.2em]
\downarrow&&\downarrow&&\downarrow&&\downarrow\\[0.45em]
\mathinner{\cdotp\cdotp\cdotp\cdotp}H^n(X;G,A_U)&\to&H^n(X;G,A)&\to&H^n(X;G,A_{X_0})&\to&
H^{n+1}(X;G,A_U)\!\to\!\mathinner{\cdotp\cdotp\cdotp\cdotp}
\end{array}
\]

\LayerTarget{tohoku:proof:transformation-first-spectral-sequence:p1}
Or il résulte de l'hypothèse sur $Y_0$ que $\mathbf H^n(G,A_U)=0$ pour $n>0$,
car c'est évidemment le cas dans l'ouvert $V$, mais c'est vrai aussi sur
$Y_0$, car en vertu

\TohokuContinuation{206}
